\begin{table}[htbp]
  \centering
  \footnotesize
  \caption{赛题口径与当前运行页面字段的对应关系}
  \label{tab:metric-interpretation}
  \begin{tabularx}{\textwidth}{>{\bfseries\raggedright\arraybackslash}p{0.15\textwidth} p{0.22\textwidth} p{0.23\textwidth} X}
    \toprule
    指标方向 & 赛题口径 & 当前页面显示 & 解释说明 \\
    \midrule
    实时性 & 平均端到端时延、TTFT 及其变化 & 边缘处理 P95、端到端 P95、页面内相对字段 & P95 与平均值不是同一统计量，端到端时延也不等同于 TTFT；页面内相对字段不在正文中延伸为独立结论。 \\
    感知与决策 & 能力保持率、准确率／召回率／F1 或任务效果 & 推理完成率、异常比例、边云配对与一致性字段 & 完成率说明处理链状态，不等同于语义质量；少量配对记录不能代表完整任务评价。 \\
    资源与通信 & 内存、上传量、带宽与计算压力 & 云端卸载率、估算减少上传、累计 Token & 估算值需要明确原始／上传字节定义，累计 Token 与最近事件窗口不可直接相除。 \\
    稳定性 & 网络波动中的基本业务保持率 & 完成率、降级比例、Defer 队列、Outbox 成功率 & 无 Outbox 事件时百分比可能是零分母，不解释为成功或失败；单次正常窗口不能替代网络扰动过程。 \\
    一致性 & 冲突率与冲突解决成功率 & 边云配对样本、一致率、冲突解决率 & 需要固定相关任务、冲突和解决的判定；无冲突时不把 0\% 直接解释为处理失败。 \\
    扩展性 & 两类差异化场景、节点与应用变化适配 & 设备／来源字段与当前登记状态 & 当前设备数量说明运行时登记，不代表架构上限；场景能力以接口与演示路径共同说明。 \\
    \bottomrule
  \end{tabularx}
\end{table}
